\chapter{Conclusions i Treball Futur}
\label{ch:conclusions}

Aquest capítol recull les conclusions finals derivades del disseny i la implementació de \textit{Web Core View}. Es realitza un balanç sobre l'assoliment dels objectius inicials, es proposen línies de millora per a futures versions i es detallen les reflexions personals sobre l'aprenentatge adquirit durant el desenvolupament del Treball de Final de Grau.

\section{Assoliment dels objectius}
Després de completar el cicle de desenvolupament i validació, es pot afirmar que s'han assolit els objectius plantejats en el Capítol 1:

\begin{itemize}
    \item \textbf{Modernització de la plataforma:} S'ha substituït amb èxit l'aplicació web \textit{legacy} dels anys 90, basada en ActiveX i tecnologies desfasades, per una interfície de gestió web moderna basada en Angular 19. La nova interfície és responsiva, ubiqüitària i accessible des de qualsevol navegador modern, sense dependència de \textit{plugins} propietaris.
    \item \textbf{Motor dinàmic (Schema-driven):} Sistema de configuració que s'adapta automàtica\-ment al maquinari de Lanaccess mitjançant JSON, eliminant la necessitat de reprogramar el \textit{frontend} per a cada \textit{firmware}.
    \item \textbf{Visualització d'alt rendiment:} Mitjançant la integració de WebRTC, s'han assolit latències inferiors als 300ms en la visualització de vídeo en directe, complint els requisits operatius de seguretat.
    \item \textbf{Gestió reactiva d'alarmes:} L'ús de WebSockets (WSS) i \textit{Angular Signals} ha permès crear una interfície que reacciona en temps real als esdeveniments del gravador sense penalitzar el rendiment del terminal de l'usuari.
    \item \textbf{Solució autònoma per a empreses petites (single-NVR):} S'ha aconseguit que l'aplicació funcioni de manera completament autònoma dins del propi videogravador, permetent a empreses petites amb un únic dispositiu gestionar-lo integralment des del navegador sense necessitat d'infraestructura centralitzada. Aquesta capacitat obre un nou segment de mercat per a Lanaccess que la plataforma anterior no cobria.
    \item \textbf{Gestió del projecte dins del pressupost i terminis:} El projecte s'ha lliurat en el termini establert de 540 hores (18 ECTS), amb un cost real de 15.028,75~\euro{}, inferior als 16.420,25~\euro{} pressupostats inicialment incloent el marge de contingència del 10\%. Les desviacions tècniques (H.265 i motor \textit{schema-driven}) es van gestionar de forma proactiva dins del marge definit, sense necessitat de finançament addicional ni retall de funcionalitats crítiques. Aquest resultat demostra que la metodologia àgil basada en Kanban i el control d'indicadors EVM (CPI/SPI) han estat eines efectives per garantir la viabilitat econòmica i temporal del projecte.
\end{itemize}

\section{Treball futur}
Tot i que el sistema és totalment funcional en el seu mode autònom \textit{single-NVR}, la naturalesa modular del projecte permet plantejar diverses línies d'evolució per a la versió 2.0:

\begin{itemize}
    \item \textbf{Analítica de vídeo amb IA:} Integrar metadades d'intel·ligència artificial directament sobre el reproductor de vídeo per dibuixar quadres de detecció d'objectes o mapes de calor en temps real.
    \item \textbf{Multi-dispositiu centralitzat (Multi-NVR):} La gestió centralitzada de múltiples gravadors des d'una única instància va formar part de l'abast inicial del projecte. Durant la fase de planificació, davant de la complexitat tècnica acumulada en la integració de H.265 i el motor \textit{schema-driven}, es va prendre la decisió estratègica de \textit{descoping} de reduir l'abast a la versió \textit{single-NVR} autònoma, prioritzant la qualitat i el lliurament dins del termini. L'evolució natural d'aquesta base consistiria a implementar el mòdul de gestió centralitzada que permeti visualitzar i configurar desenes de gravadors Lanaccess distribuïts geogràficament des d'una única instància de l'aplicació, atenent el segment corporatiu de grans instal·lacions.
    \item \textbf{Suport d'enregistrament avançat:} Ampliar les eines d'edició i exportació de vídeo directament des del navegador, permetent l'aplicació de marques d'aigua digitals o signatures legals en les descàrregues.
    \item \textbf{Aplicació Mòbil Nativa:} Utilitzar tecnologies com \textit{Ionic} o \textit{Capacitor} per reutilitzar el codi d'Angular 19 i oferir una versió nativa a les botigues d'aplicacions (App Store / Play Store).
\end{itemize}

\section{Conclusions personals}
La realització d'aquest TFG ha estat una experiència d'aprenentatge integral que ha anat molt més enllà de la simple programació. 

Tècnicament, m'ha permès dominar les darreres funcionalitats d'Angular 19, especialment el canvi de paradigma que suposen els \textit{Signals} per a la gestió de l'estat reactiu. Així mateix, el repte de treballar amb protocols de vídeo com WebRTC i HLS m'ha proporcionat una visió profunda sobre la gestió de fluxos multimèdia en entorns de xarxa complexos.

Quant a la gestió, el projecte m'ha ensenyat la importància de l'arquitectura \textit{schema-driven}. Desenvolupar un software que no sàpiga exactament què ha de configurar fins que rep les dades del servidor requereix un nivell d'abstracció i rigor en el disseny que ha posat a prova els meus coneixements d'arquitectura de software i principis SOLID.

Finalment, el fet de desenvolupar el projecte per a una empresa real com Lanaccess m'ha permès entendre les necessitats de mantenibilitat i escalabilitat que exigeix el mercat industrial, on el software ha de ser, per sobre de tot, robust i eficient. L'èxit de la visualització de baixa latència ha estat possible gràcies a l'estreta coordinació entre el frontend i el microservei \textbf{la-video-engine}, demostrant la viabilitat de WebRTC en entorns industrials.